By Corollary~\ref{cor:Aprime},
\[ C_n=2\,(4/\pi)^{m/2}\,\Gamma(2+m/2)\prod_{\chi}|L(1,\chi)|\big/\log(2+\sqrt5)^m , \]
the product over the $m$ even Dirichlet characters of conductor exactly $q=2^{n+2}$ (Proposition~\ref{prop:D}). Each such $\chi$ is primitive and even, $2\mid q$ and $q\ge16\ge4$, so \cite{Ram04} Corollary~1 with $h=1$ and $k=2$ gives $0\le|L(1,\chi)|\le U(q)$. A product of $m$ real numbers in $[0,U(q)]$ is at most $U(q)^m$; hence
\[ C_n\le2\,(4/\pi)^{m/2}\,\Gamma(2+m/2)\,U(q)^m\big/\log(2+\sqrt5)^m=\widehat C_n . \] If $\ell^{d_n(\ell)}>\widehat C_n$ then $\ell^{d_n(\ell)}>C_n$, and Theorem~\ref{thm:A} gives $\ell\nmid k_n$. The discrete part (a product of $m$ factors in $[0,U]$ is $\le U^m$, and $C\le\widehat C<\ell^d$ feeds $C<\ell^d$ into the discrete core \path{TheoremAInputs.theoremA_of_inputs}) is the Lean theorems \path{WeberHatC.prod_le_pow_card_of_le}, \path{WeberHatC.C_le_hatC} and \path{WeberHatC.theoremA_of_hat}; the analytic input $|L(1,\chi)|\le U(q)$ is a literature hypothesis and is not formalised.
